\documentclass[a4paper]{article}     
\usepackage{amsfonts}
\usepackage{amsmath}
\usepackage{amssymb}
\usepackage{graphicx}
\usepackage{color}

\newcommand{\Q}{\mathbb{Q}}
\newcommand{\R}{\mathbb{R}}
\newcommand{\llim}{\lim\limits}
\newcommand{\dint}{\displaystyle\int}

\begin{document}

\noindent \large Auteur: Ilyas Sefiani \\
\noindent \large Studierichting: Informatica\\
\noindent \large Oefeningenreeks 5B nr. 5.10\\

\medskip

\normalsize

$\dint_{0}^{\tfrac{\pi}{2}} \dfrac{\sin(x)dx}{\sin(x)+\cos(x)}$\\

Primitieve vinden:\\

$$t = tan\left(\dfrac{x}{2}\right)$$ $$dx = \dfrac{2}{1+t^2}dt$$\\

$\displaystyle\int \dfrac{\dfrac{2t}{1+t^2}}{\dfrac{-t^2+2t+1}{1+t^2}}\dfrac{2dt}{1+t^2}$\\

$=\displaystyle\int \dfrac{4t}{(-t^2+2t+1)(1+t^2)}dt$\\

$=\displaystyle\int -\dfrac{4t}{(t-(1+\sqrt{2}))(t-(1-\sqrt{2}))(1+t^2)}dt$\\

Stelling Fuss (3)\\

$-\dfrac{4t}{(t-(1+\sqrt{2}))(t-(1-\sqrt{2}))(1+t^2)} = \dfrac{A}{t-(1+\sqrt{2})} + \dfrac{B}{t-(1-\sqrt{2})} + \dfrac{Ct+D}{(1+t^2)}$\\

Voor $t = 1+\sqrt{2}$:\\

$A = \dfrac{-4t}{(t-(1-\sqrt{2}))(1+t^2)} = \dfrac{-4(1+\sqrt{2})}{(2\sqrt{2})(4+2\sqrt{2})} = \dfrac{-4(1+\sqrt{2})}{(8\sqrt{2}+8)} = \dfrac{-4}{8} = -\dfrac{1}{2}$\\

Voor $t = 1-\sqrt{2}$:\\

$B = \dfrac{-4t}{(t-(1+\sqrt{2}))(1+t^2)}=\dfrac{-4(1-\sqrt{2})}{(-2\sqrt{2})(4-2\sqrt{2})}=\dfrac{-4(1-\sqrt{2})}{(-8\sqrt{2}+8)}=-\dfrac{4}{8}$\\

$= -\dfrac{1}{2}$\\

Voor $t=i$:\

$Ci+D = \dfrac{-4t}{(t-(1+\sqrt{2}))(t-(1-\sqrt{2}))}=\dfrac{-4i}{(i-(1+\sqrt{2}))(i-(1-\sqrt{2}))}$\\

$=\dfrac{-4i}{(i-1-\sqrt{2}))(i-1+\sqrt{2}))}=\dfrac{-4i}{(-2i-2)}=\dfrac{2i}{i+1}$\\

$=\dfrac{2i(i-1)}{(i+1)(i-1)}=\dfrac{-2-2i}{-2}=i+1$\\

$(C,D) = (1,1)$ dus $C=1$ en $D=1$\\

$\dfrac{4t}{(-t^2+2t+1)(1+t^2)} = \dfrac{-1}{2(t-(1+\sqrt{2}))} + \dfrac{-1}{2(t-(1-\sqrt{2}))} + \dfrac{t+1}{(1+t^2)}$\\

Integralen berkenen:\\

$-\dfrac{1}{2}\displaystyle\int \dfrac{dt}{t-(1+\sqrt{2})}$\\

$$u=t-(1+\sqrt{2})$$ $$du=dt$$\\

$=-\dfrac{1}{2}\displaystyle\int \dfrac{du}{u} = -\dfrac{1}{2}\ln(u) + C = -\dfrac{1}{2}\ln|t-(1+\sqrt{2})| + C$\\


$-\dfrac{1}{2}\displaystyle\int \dfrac{dt}{t-(1-\sqrt{2})}$\\

$$u=t-(1-\sqrt{2})$$ $$du=dt$$\\

$=-\dfrac{1}{2}\displaystyle\int \dfrac{du}{u} = -\dfrac{1}{2}\ln(u) + C = -\dfrac{1}{2}\ln|t-(1-\sqrt{2})| + C$\\

$\displaystyle\int \dfrac{t+1}{t^2+1} = \displaystyle\int \dfrac{t}{t^2+1}dt + \displaystyle\int \dfrac{1}{t^2+1}dt$\\

$$u=t^2+1$$ $$du=2tdt$$\\

$=\dfrac{1}{2}\displaystyle\int \dfrac{du}{u} + \operatorname{Bgtan}{t} = \dfrac{1}{2}\ln|u| + \operatorname{Bgtan}{t} + C= \dfrac{1}{2}\ln|t^2+1| + \operatorname{Bgtan}{t} + C$\\

Samengesteld:\\

$-\dfrac{1}{2}\ln|(t-(1+\sqrt{2}))(| - \dfrac{1}{2}\ln|t-(1-\sqrt{2})| + \dfrac{1}{2}\ln|t^2+1| + \operatorname{Bgtan}{t} + C$\\

$= \dfrac{1}{2}\ln\left|\dfrac{t^2+1}{(t-(1+\sqrt{2}))(t-(1-\sqrt{2}))}\right| + \operatorname{Bgtan}{t} + C$\\

$t = \tan\left(\dfrac{x}{2}\right)$:\\

We weten $\sin(x)+\cos(x) = \dfrac{-t^2+2t+1}{1+t^2}$ dus $\dfrac{t^2+1}{t^2-2t-1}=-\dfrac{1}{\sin(x)+\cos(x)}$:\\

En $\operatorname{Bgtan}{\left(\tan\left(\dfrac{x}{2}\right)\right)} = \dfrac{x}{2}$:\\

Dus $-\dfrac{1}{2}\ln|t-(1+\sqrt{2})| - \dfrac{1}{2}\ln|t-(1-\sqrt{2})| + \dfrac{1}{2}\ln|t^2+1| + \dfrac{x}{2} + C$\\

$= \dfrac{1}{2}\ln\left|-\dfrac{1}{\sin(x)+\cos(x)}\right| + \operatorname{Bgtan}{t} + C = \dfrac{1}{2}\ln\left|-\dfrac{sin^2(x)+cos^2(x)}{\sin(x)+\cos(x)}\right| + \dfrac{x}{2} + C$\\

$= -\dfrac{1}{2}\ln|\sin(x)+\cos(x)| + \dfrac{x}{2} + C$\\

Grenzen invullen:\\

$-\dfrac{1}{2}\ln|1| + \dfrac{\pi}{4} + \dfrac{1}{2}\ln|1| - \dfrac{0}{2} = \dfrac{\pi}{4}$\\

De integraal is convergent!\\

\bigskip

\emph{Opmerking WP: Dat is juist, maar er is een veel kortere methode die ik u niet wil onthouden. Stel
$$I = \dint_{0}^{\tfrac{\pi}{2}} \dfrac{\sin(x)dx}{\sin(x)+\cos(x)}$$
Stel $t=\dfrac{\pi}{2}-x$, dan is
$$I = -\dint_{\tfrac{\pi}{2}}^0 \dfrac{\cos(t)dt}{\cos(t)+\sin(t)}$$
Tellen we de twee integralen op, dan is 
$$2I = \dint_{0}^{\tfrac{\pi}{2}} \dfrac{(\sin(x)+\cos(x))dx}{\sin(x)+\cos(x)} = \dint_{0}^{\tfrac{\pi}{2}} dx = \dfrac{\pi}{2}$$
dus $I=\dfrac{\pi}{4}$}

\begin{thebibliography}{999}
\end{thebibliography}
\end{document} 